\documentclass[letterpaper]{article}
\usepackage[submission]{aaai2027}
% The serif, sans-serif, and monospaced fonts are loaded automatically by
% aaai2027.sty. Do not add font packages.
\usepackage[hyphens]{url}
\usepackage{graphicx}
\urlstyle{rm}
\def\UrlFont{\rm}
\usepackage{natbib}
\usepackage{caption}
\frenchspacing
\usepackage{booktabs}
\usepackage{amsmath}
\usepackage{amssymb}

\pdfinfo{
/TemplateVersion (2027.1)
}

\setcounter{secnumdepth}{1}

\newcommand{\R}{\mathbb{R}}
\newcommand{\Attn}{\mathrm{Attn}}
\newcommand{\FFN}{\mathrm{FFN}}
\newcommand{\LN}{\mathrm{LN}}
\newcommand{\BA}{B^{A}}
\newcommand{\BF}{B^{F}}
\newcommand{\ca}{c^{A}}
\newcommand{\cf}{c^{F}}

\title{A Write-Economy View of Transformer Residual Blocks}

\author{
    Anonymous Submission
}
\affiliations{
}

\begin{document}

\maketitle

\begin{abstract}
Residual connections are usually treated as optimization devices or shared
state pathways. We study a more structural role: a residual block writes
learned directions into a persistent state with coefficients generated from
that same state. Transformer residual streams provide a particularly rich
instance because Attention heads, feed-forward networks (FFNs), embeddings, and
unembeddings all communicate through the same stream. This view, while useful,
does not distinguish two architecturally different roles: directly writing
learned directions into the residual stream and merely modulating the
coefficients of another module's writes. We use a basis/coefficient view
of Transformer blocks in which Attention output projections and FFN down
projections define distinct learned write-basis families, while their
coefficients are generated dynamically from the mixed residual history. Under
this view, the standard Transformer block is not only a composition of
Attention and FFN computations; it is a dual-write system in which both modules
retain direct residual-stream write access.

We test this interpretation through controlled decoder-only language-modeling
experiments on enwik8. In a five-seed, parameter-matched topology sweep,
standard Attention-then-FFN blocks obtain the lowest test loss. Reversing order
while preserving both direct write families causes only a small paired
degradation ($+0.0145$ test NLL), whereas removing one direct write family is
substantially more damaging: coefficient-only block variants degrade by
$+0.0274$ to $+0.0560$, and carry variants remain weaker despite giving the
removed module additional influence over coefficient generation. A
complementary output-projection absorption control is consistent with a
distinction between a projection as an internal parameterization and as a
residual write outlet. These results do not by themselves isolate direct write
access from all normalization and optimization-path effects, but they support a
write-economy hypothesis: residual streams are shared spaces where model
quality depends on which heterogeneous modules can write directly and how their
dynamic coefficients are coupled.
\end{abstract}

\section{Introduction}

Residual connections appear throughout modern deep learning, from convolutional
networks to language models \citep{he2016identity,vaswani2017attention}. They
are usually explained as identity paths that stabilize optimization, preserve
signals, or let layers refine a shared representation. We ask a complementary
functional question: when a residual block updates its state, what learned
directions does it write through, and which computation generates the
coefficients of those directions?

The Transformer residual stream is a particularly revealing case. Token
embeddings, Attention heads, feed-forward networks (FFNs), normalization
layers, and the unembedding all read from or write to the same state
\citep{vaswani2017attention,elhage2021framework}. It is therefore natural to
view the stream as a representation space or memory substrate: a state whose
geometry, features, or linear probes reveal what the model stores at a given
layer \citep{nostalgebraist2020logit,belrose2023tuned}. This view is necessary
but not sufficient. A memory substrate is only useful through its access
mechanism. In Transformers, the corresponding architectural question is: which
modules are allowed to directly write their own learned directions into the
residual stream?

We argue that this distinction matters. A sublayer may influence the residual
stream in two different ways. It may directly write an update through its own
output basis, or it may only shape the coefficients used by a later sublayer.
These roles are conflated in the standard description of Transformer blocks,
where Attention and FFN are often treated as two nonlinear transformations
connected by residual paths. In a standard pre-normalization block, however,
both sublayers are direct authors of the residual stream: Attention writes
through its output projection, and the FFN writes through its down projection.
In contrast, block-composed variants such as $h+\FFN(\Attn(h))$ allow
Attention to affect FFN coefficients but remove Attention's direct write
outlet.

We formalize this distinction through a basis/coefficient view of residual
updates. The columns of an Attention output projection form an Attention
write-basis family, while the columns of an FFN down projection form an FFN
write-basis family. The corresponding coefficients are not fixed parameters:
Attention coefficients combine context-dependent routing weights with value
coordinates, and FFN coefficients are nonlinear activations generated from the
current residual state. A standard block therefore produces updates of the form
\begin{equation}
\Delta H_l = \BA_l \ca_l(H_l) +
             \BF_l \cf_l(H_l, \BA_l \ca_l(H_l)),
\label{eq:dual-write}
\end{equation}
where both basis families write directly and both coefficient functions are
coupled through the mixed residual history.

This view makes a concrete prediction: architectures that preserve both direct
write families should remain closer to the standard Transformer than
architectures that allow one module only to modulate another module's
coefficients. We test this prediction by comparing standard, parallel,
order-reversed, block-composed, and carry variants under matched
language-modeling settings. The carry variants are especially diagnostic: they
pass additional intermediate signals across layers, giving the removed module a
stronger role in coefficient generation, but they still do not restore its
direct write basis.

Our contributions are threefold. First, we formulate Transformer residual
updates in terms of architecturally fixed write outlets and dynamic
coefficients, separating direct residual write access from coefficient
modulation. Second, we derive the Attention and FFN write-outlet families for
standard decoder-only Transformers and show how their coefficients depend on
residual history. Third, we provide controlled language-modeling experiments
showing that the standard dual-residual topology is consistently stronger than
the tested order-reversed, parallel, block-composed, and carry variants under a
matched five-seed setting. We treat pretrained-model attribution and
counterfactual write patching as necessary follow-up evidence rather than as
completed contributions in the present draft.

\section{Residual Streams as Write Economies}

A generic residual block has the form
\begin{equation}
x_{l+1}=x_l+g_l(x_l).
\end{equation}
When the last operation in $g_l$ is a learned linear map, the update can often
be written as
\begin{equation}
g_l(x_l)=B_l c_l(x_l),
\label{eq:generic-write}
\end{equation}
where $B_l$ is a learned family of write directions and $c_l(x_l)$ is a
state-dependent coefficient vector. This covers fully connected residual MLP
blocks directly. In convolutional residual blocks, the same idea applies with
spatially shared convolutional output channels acting as write dictionaries
over local neighborhoods. Transformer blocks are a more structured case: within
one residual stream, multiple heterogeneous submodules own different output
dictionaries and generate coefficients in different ways.

Let $H_l=[h_{l,1},\ldots,h_{l,T}]\in\R^{d\times T}$ denote the residual stream
at layer $l$, with token vectors written as columns. A standard pre-LayerNorm
Transformer block can be written as
\begin{align}
U_l &= H_l + A_l(\LN(H_l)),\\
H_{l+1} &= U_l + F_l(\LN(U_l)).
\end{align}
Equivalently,
\begin{equation}
H_{l+1}=H_l+\Delta_l^A+\Delta_l^F.
\end{equation}
The usual residual-path view emphasizes the identity term $H_l$. Our focus is
the authorship of the update terms: which learned direction families can write
directly into the shared stream.

We call a learned matrix column family a \emph{write basis} when it is the final
linear outlet through which a submodule writes into the residual stream. This
does not require linear independence; the basis may be overcomplete,
correlated, and redundant. A module has \emph{direct write access} if its own
write basis appears in the residual update. A module has only
\emph{coefficient access} if it can affect another module's coefficient
generator but cannot write its own final directions.

This definition is architectural rather than coordinate-free. In typical
Transformers, both $W_O$ and the FFN down projection may have column spans equal
to the full residual space, so our claim is not that each module owns a
disjoint reachable subspace. Nor is an individual column intrinsically
identified: for an invertible change of coordinates $R$, $Bc=(BR)(R^{-1}c)$.
The object we study is the parameterized outlet fixed by the architecture and
training dynamics. Direct write access therefore means an independent learned
map from a module's internal coordinates into the residual update, not exclusive
ownership of a linear subspace.

\section{Attention and FFN Write Bases}

\paragraph{Attention.}
For one Attention head $r$, let
$q_i^r=W_Q^r x_i$, $k_j^r=W_K^r x_j$, and $v_j^r=W_V^r x_j$, where
$x_i$ is the normalized residual vector at position $i$. The causal routing
weights are
\begin{equation}
\alpha_{ij}^r =
\frac{\exp((q_i^r)^\top k_j^r/\sqrt{d_h})}
{\sum_{m\leq i}\exp((q_i^r)^\top k_m^r/\sqrt{d_h})}.
\end{equation}
Writing the output projection as $W_O=[O_1,\ldots,O_h]$ and the columns of
$O_r$ as $o_{r,a}$, the Attention residual update at token $i$ is
\begin{equation}
\Delta_{l,i}^{A}
=\sum_{r=1}^{h}\sum_{a=1}^{d_h}
\left(\sum_{j\leq i}\alpha_{ij}^r v_{j,a}^r\right)o_{r,a}.
\label{eq:attn-basis}
\end{equation}
Thus the Attention write basis consists of the columns of $W_O$, while the
coefficients combine routing weights and value coordinates.

\paragraph{FFN.}
For a standard FFN,
\begin{equation}
F_l(z)=W_l^2\phi(W_l^1z+b_l^1)+b_l^2.
\end{equation}
Let $v_{l,k}$ be the $k$-th column of $W_l^2$. Ignoring the output bias,
\begin{equation}
F_l(z)=\sum_k
\phi(\langle u_{l,k},z\rangle+b_{l,k})v_{l,k}.
\label{eq:ffn-basis}
\end{equation}
The FFN write basis is therefore the column family of the down projection, and
the coefficient of each direction is a nonlinear detector response computed
from the current residual state.

\paragraph{Coupled coefficients.}
The write bases are static learned outlets, but the coefficients are functions
of residual history. In a standard Attention-then-FFN block, the FFN
coefficients depend on the Attention write made earlier in the same block.
Therefore the standard block is not the sum of two independent
basis/coefficient pairs; it is a coupled dual-write system.

\section{Architectures That Separate Writes and Coefficients}

Table~\ref{tab:variants} summarizes the variants used to separate direct write
access from ordering and coefficient coupling. The key contrast is between
architectures that preserve both direct write families and architectures where
one module only modulates coefficients before the other module writes.

\begin{table*}[t]
\centering
\caption{Architecture variants separate direct write access, ordering, and
coefficient coupling. A checkmark indicates that the module writes its own
output basis directly into the residual stream.}
\label{tab:variants}
\begin{tabular}{lllll}
\toprule
Variant & Update form & Attn write & FFN write & Main contrast \\
\midrule
standard & $h+\Attn(h)+\FFN(h+\Attn(h))$ & yes & yes & full dual-write block \\
standard\_fa & $h+\FFN(h)+\Attn(h+\FFN(h))$ & yes & yes & order reversed \\
parallel & $h+\Attn(h)+\FFN(h)$ & yes & yes & same-layer AF coupling removed \\
block\_af & $h+\FFN(\Attn(h))$ & no & yes & Attention coefficient-only \\
block\_fa & $h+\Attn(\FFN(h))$ & yes & no & FFN coefficient-only \\
block\_af\_carry & $h+\FFN(\Attn(h)+\Attn(h_{-1}))$ & no & yes & Attention carry only \\
block\_fa\_carry & $h+\Attn(\FFN(h)+\FFN(h_{-1}))$ & yes & no & FFN carry only \\
\bottomrule
\end{tabular}
\end{table*}

\section{Controlled Language-Model Experiments}

\paragraph{Controlled ablation protocol.}
We evaluate decoder-only language models on enwik8 character-level language
modeling. All compared variants are implemented in the same codebase, trained
from scratch on identical train/validation/test splits, and evaluated with the
same tokenizer, context length, model width, depth, dropout, and optimization
schedule unless explicitly noted. Checkpoints are selected by validation loss
and evaluated once on the held-out test split. Unless otherwise stated, models
use 8 layers, 8 heads, hidden size 512, context length 512, and batch size 256.
The main sweep uses the Muon optimizer and five random seeds. All seven main
variants have 25.6M trainable parameters. We report validation and test loss;
lower is better. Paired deltas compare each variant with the standard block at
the same seed and use two-sided 95\% $t$ intervals over the five paired
differences. The carry variants share weights between the current-state and
previous-state branches, so their parameter counts remain matched to the
corresponding block variants. They do, however, invoke one submodule twice per
block. Thus their underperformance is not explained by receiving less
computation than the standard dual-write block.

\paragraph{Basis-carry comparison.}
Table~\ref{tab:basis-carry} isolates the key carry comparison from the main
five-seed sweep. Reversing the standard AF order while keeping both direct
write families causes a modest degradation. In contrast, carry variants remain
much weaker even though they give the removed module additional coefficient
influence through a previous middle signal.

\begin{table*}[t]
\centering
\caption{Basis-carry result on enwik8 from the five-seed Muon sweep. Losing one
direct write family is more damaging than reversing sublayer order while
preserving both direct write families. Test losses are mean $\pm$ sample
standard deviation; deltas are paired against the standard block.}
\label{tab:basis-carry}
\begin{tabular}{llllrr}
\toprule
Variant & Attn write & FFN write & Carry signal & Test loss $\downarrow$ & Delta $\downarrow$ \\
\midrule
standard & yes & yes & no & $0.8358 \pm 0.0027$ & $0.0000$ \\
standard\_fa & yes & yes & no & $0.8503 \pm 0.0061$ & $+0.0145$ \\
block\_af\_carry & no & yes & yes & $0.8617 \pm 0.0026$ & $+0.0259$ \\
block\_fa\_carry & yes & no & yes & $0.8780 \pm 0.0070$ & $+0.0422$ \\
\bottomrule
\end{tabular}
\end{table*}

\paragraph{Five-seed topology sweep.}
Table~\ref{tab:topology} reports the complete Muon-optimized topology sweep.
All non-standard variants are worse than the standard block in all five paired
seeds. The ordering matches the write-economy prediction. The closest variant
is \texttt{standard\_fa}, which reverses order but preserves both direct write
families. The parallel variant also preserves both direct writes but removes
same-layer AF coefficient coupling, producing an intermediate gap. The
block-composed variants, which turn one module into coefficient-only access,
are much weaker. The asymmetry is also informative: in this setting, denying
the FFN its direct write basis (\texttt{block\_fa}) is more damaging than
denying Attention its direct write basis (\texttt{block\_af}).

\begin{table*}[t]
\centering
\caption{Five-seed topology sweep on enwik8 with 25.6M-parameter models. Test
losses are mean $\pm$ sample standard deviation. Paired deltas and confidence
intervals compare each variant to \texttt{standard} at the same seed. Every
non-standard variant is worse than \texttt{standard} in all five paired seeds.}
\label{tab:topology}
\begin{tabular}{lrrrr}
\toprule
Variant & Test NLL $\downarrow$ & Test bpc $\downarrow$ & $\Delta$ test NLL $\downarrow$ & $\Delta$ \% \\
\midrule
standard & $0.8358 \pm 0.0027$ & $1.2058 \pm 0.0039$ & -- & -- \\
standard\_fa & $0.8503 \pm 0.0061$ & $1.2268 \pm 0.0088$ & $+0.0145$ $[0.0086,0.0204]$ & $+1.74$ \\
parallel & $0.8561 \pm 0.0024$ & $1.2350 \pm 0.0035$ & $+0.0202$ $[0.0186,0.0219]$ & $+2.42$ \\
block\_af & $0.8632 \pm 0.0030$ & $1.2453 \pm 0.0043$ & $+0.0274$ $[0.0251,0.0297]$ & $+3.28$ \\
block\_af\_carry & $0.8617 \pm 0.0026$ & $1.2432 \pm 0.0038$ & $+0.0259$ $[0.0248,0.0270]$ & $+3.10$ \\
block\_fa & $0.8918 \pm 0.0055$ & $1.2866 \pm 0.0079$ & $+0.0560$ $[0.0502,0.0617]$ & $+6.70$ \\
block\_fa\_carry & $0.8780 \pm 0.0070$ & $1.2667 \pm 0.0101$ & $+0.0422$ $[0.0330,0.0514]$ & $+5.05$ \\
\bottomrule
\end{tabular}
\end{table*}

\paragraph{Output-projection absorption control.}
In a block-AF topology without an intervening nonlinearity or normalization
between Attention and FFN, the Attention output projection $W_O$ can be absorbed
into the FFN input projection. Empirically, removing $W_O$ in this setting is
quality-neutral or slightly beneficial, while the original block-AF with a
middle normalization remains stronger. This control supports the distinction
between a projection as an internal parameterization and a projection as a
direct residual write outlet, but it should be interpreted cautiously: the
experiment mainly checks an algebraic redundancy of consecutive linear maps and
is not, by itself, a causal proof of the write-access hypothesis.

\section{Pretrained Open-Model Diagnostics}

The controlled experiments isolate architecture mechanisms, but the same
write-outlet decomposition also maps onto mainstream pretrained decoder-only
models. In GPT-NeoX/Pythia models, Attention output outlets correspond to
Attention dense output projections and FFN outlets to dense-to-hidden output
projections. In GPT-2, both appear as output \texttt{c\_proj} modules. In
Llama/Qwen-style models, they correspond to \texttt{o\_proj} and
\texttt{down\_proj}. A complete version of this section should go beyond this
inventory: it should attribute token logits to module-level writes and then
test those attributions with clean/corrupt counterfactual write patching and
matched controls. We leave the TODOs explicit in the current draft to avoid
claiming evidence not yet run.

% TODO: Insert open-model inventory table after running inspect_model_basis.py.
% TODO: Insert module-level logit attribution and counterfactual write patching.

\section{Related Work}

\paragraph{Residual streams, memory, and Transformer circuits.}
Transformer Circuits frames the residual stream as a shared communication space
and decomposes Attention through QK and OV circuits
\citep{elhage2021framework}. Recent architectural work further treats the
residual stream as a scalable memory substrate, for example by replacing vector
residual states with matrix-valued residual memories \citep{mak2025residual}.
Our work is complementary but more fine-grained. We do not ask whether the
stream can store information or whether its capacity can be enlarged; we ask
how the architecture allocates direct write access to that stream. In this
sense, the write-economy view refines the memory-bus view by separating the
storage medium from the write bases and coefficient generators that use it.

\paragraph{Residual optimization, normalization, and scaling.}
Residual connections are known to ease optimization in deep networks
\citep{he2016identity}. In Transformers, a large body of work modifies
normalization placement, residual scaling, initialization, or branch gating to
stabilize deep training. Pre-LN analysis explains why moving normalization
inside residual blocks improves gradient behavior \citep{xiong2020layer};
ScaleNorm/FixNorm and T-Fixup-style initialization reduce dependence on warmup
or normalization \citep{nguyen2019transformers,huang2020improving}; Admin
links Transformer training difficulty to residual-branch amplification and
uses adaptive initialization to stabilize early training
\citep{liu2020understanding}; ReZero and LayerScale gate residual branches
\citep{bachlechner2021rezero,touvron2021going};
NormFormer adds extra normalization and head-wise scaling around sublayer
outputs \citep{shleifer2021normformer}; and DeepNorm modifies residual
normalization and initialization to scale Transformers to extreme depth
\citep{wang2022deepnet}. ResiDual further combines Pre-LN and Post-LN residual
routes to address both gradient and representation-collapse issues
\citep{xie2023residual}. These works primarily regulate the magnitude,
conditioning, or gradient behavior of residual updates. Our analysis is
complementary: it asks which learned basis family owns direct write access once
an update is allowed to enter the stream.

\paragraph{Residual connections and representation collapse.}
Another line of analysis studies residual connections as a source of
representational expressivity rather than only trainability. Pure self-attention
networks without skip connections or MLPs have been shown to converge toward
token-uniform low-rank representations, while skip connections and MLPs prevent
this degeneration \citep{dong2021attention}. This supports the broader view
that residual pathways and FFNs play functional roles in Transformer
computation. Our contribution is more specific: conditioned on using residual
updates, we ask which heterogeneous output dictionaries are allowed to write
directly into the shared state.

\paragraph{Residuals in routing and cross-layer reads.}
Not all residual variants operate on the final hidden-state update. RealFormer
adds residual connections over Attention scores and reports more stable and
sparser Attention \citep{he2021realformer}. Deep Transformer variants also pass
combinations of previous layers to later layers \citep{wang2019learning}. In
our terminology, such methods modify coefficient or read history: they change
the states from which routing and activations are generated. This differs from
restoring a missing final write basis. A module can receive richer history or
stronger routing residuals while still being denied direct residual-stream
write access.

\paragraph{Parallel and alternative Transformer blocks.}
Parallel Attention/FFN designs appear in large language models and have been
studied as alternatives to the serial block
\citep{black2022gptneox,chowdhery2023palm,sonkar2023paf}. Our framework
explains why parallel variants can remain competitive: they preserve both
direct write families while weakening same-layer coefficient coupling.

\paragraph{Parameter sharing and conditional experts.}
Universal Transformers and ALBERT-style models reuse transformation parameters
across depth to improve parameter efficiency or add a recurrent inductive bias
\citep{dehghani2018universal,lan2019albert}. In the opposite direction,
sparsely gated mixture-of-experts layers, GShard, and Switch Transformers
increase parameter capacity through conditional expert routing
\citep{shazeer2017outrageously,lepikhin2020gshard,fedus2021switch}. These
architectures suggest future tests for the write-economy view: parameter
sharing may reduce the number of independent layer-indexed write outlets,
whereas MoE may conditionally expand the set of FFN output outlets. We do not
use these claims as evidence in this paper because parameter sharing and expert
routing also change capacity, optimization, and computation.

\paragraph{FFN memories, attribution, and learned dictionaries.}
FFNs have been interpreted as key-value memories or learned feature
dictionaries \citep{geva2021transformer}. Sparse autoencoders and activation
steering study learned or discovered directions in residual space
\citep{bricken2023monosemanticity,turner2023activation}. Our architectural
basis view is complementary: it identifies write dictionaries already present
in the model's parameterization.

\section{Discussion and Limitations}

The results suggest that residual connections should not be viewed only as
optimization aids. More generally, residual blocks allocate write rights into a
persistent computational state. In a simple residual MLP or convolutional
network, this allocation may be dominated by a single output dictionary per
block. Transformers expose a richer version of the same phenomenon: multiple
heterogeneous modules write to the same residual stream, and model quality
depends on which of them retain direct write outlets. For the model scales
studied here, Attention and FFN benefit from retaining separate, heterogeneous
direct write outlets.

This perspective also gives a common language for residual enhancements and
alternatives. Branch gates, residual scaling rules, and normalization changes
control the amplitude and conditioning of writes. Residual attention and
cross-layer aggregation enrich coefficient generation and read history. Dual
residual paths alter how normalized and unnormalized states coexist. All of
these designs can be described inside the same write economy by specifying
three objects: the state read by a coefficient generator, the basis family that
receives direct write access, and the scale or gate applied to the resulting
write. This suggests a broader use of the framework beyond the specific block
variants tested here.

The same language suggests hypotheses about capacity interventions such as
looped Transformers and mixture-of-experts layers. Looped or parameter-shared
Transformers reuse the same block parameters across depth, which may reduce the
number of independent layer-indexed write outlets. MoE layers may move in the
opposite direction by routing tokens to expert-specific FFN output outlets.
These are useful design hypotheses, but they remain outside the evidence
boundary of the present controlled sweep.

Several limitations remain. First, although the formalism applies to residual
networks beyond Transformers, our controlled experiments target decoder-only
Transformer language models because their Attention/FFN structure separates
multiple write families cleanly. We do not claim to have experimentally
characterized residual write economies in convolutional networks or residual
MLPs. Second, the controlled experiments use relatively small models on enwik8,
so the paper does not claim a complete scaling law for large language models or
a modern subword language-modeling benchmark. Third, the current topology sweep
does not uniquely isolate direct write access from all changes in normalization
placement, composition depth, Jacobian paths, and optimization geometry. A
stronger causal test would vary the rank or capacity of a direct residual
outlet while keeping coefficient pathways and compute-matched controls fixed.
Fourth, the term ``basis'' means an overcomplete learned write dictionary, not
a linearly independent or identifiable basis. Finally, pretrained open-model
attribution and intervention experiments are necessary to strengthen the
connection between the controlled architecture results and deployed language
models.

\section{Conclusion}

We used a basis/coefficient view of residual updates to distinguish direct
write access from coefficient-only modulation. In the controlled Transformer
language models studied here, the standard dual-residual topology is
consistently stronger than the tested alternatives that reverse order, remove
same-layer coupling, or make one module coefficient-only. These results are
consistent with the hypothesis that independent Attention and FFN residual
outlets are useful architectural resources, while also identifying the next
experiments needed to isolate that hypothesis more cleanly.

\bibliography{references}

\end{document}
